\begin{lemma}[Moser's local boundedness estimate]
Let $n\ge2$ and $Q>0$, and suppose the Sobolev--Poincar\'e inequality on balls
holds in dimension $n$ (\noderef{CitedSobolevPoincareBall}).  Then there is a
constant $C=C(n,Q)>0$ with the following property.  Let $B_R=B_R(x_0)\subseteq\R^n$
with $R>0$, let $v\colon\R^n\to\R$ be measurable with $v\ge0$ everywhere and
with $v^2$ integrable on $B_R$, and for every $\theta>1$ let
$g_\theta\colon\R^n\to\R^n$ be a vector field such that
\begin{enumerate}
\item $v^\theta\in W^{1,2}_{\mathrm{loc}}(B_R)$ with weak gradient $g_\theta$
  (\noderef{MemW1pLocOn}); and
\item for every $\theta>1$ and every $\zeta\in C_0^\infty(\R^n)$ with
  $\operatorname{supp}\zeta$ compact and contained in $B_R$,
  \[
  \int_{B_R}|g_\theta|^2\zeta^2\,dx
  \ \le\ Q\int_{B_R}\bigl(v^\theta\bigr)^2|\nabla\zeta|^2\,dx ,
  \]
  with one and the same $Q$ for every $\theta$.
\end{enumerate}
Then
\[
v(x)\ \le\ C\Bigl(R^{-n}\int_{B_R}v^2\,dy\Bigr)^{\frac12}
\qquad\text{for almost every }x\in B_{R/2}(x_0).
\]

\textbf{What the hypothesis says and why its uniformity is the point.}  The
inequality in (2) is a Caccioppoli, or reverse-Poincar\'e, inequality for the
single power $v^\theta$; the content of the hypothesis is that one constant $Q$
serves \emph{every} exponent $\theta>1$.  A $\theta$-dependent constant would
not do: the conclusion is reached by iterating (2) along an exponent sequence
$\theta_j\to\infty$ and multiplying the constants produced at each step, and
that product converges only because the constant does not grow with $\theta$.
This is also why the conclusion's constant depends on $n$ and $Q$ alone.

\textbf{On the two regularity hypotheses.}  Measurability of $v$ and
integrability of $v^2$ on $B_R$ are hypotheses the statement must carry even
though they go unmentioned where this estimate is used.  Measurability is what
makes the essential supremum of $v$, and hence the conclusion, meaningful.
Integrability of $v^2$ on $B_R$ is needed because $B_R$ is not compactly
contained in itself, so hypothesis (1) --- which controls $v^\theta$ only on
compact subsets of $B_R$ --- does not supply it; without it the right-hand side
would be the integral of a non-integrable function and the inequality would
assert $v\le0$.  Both are automatic at the point of use, where $v$ is a
continuous function and the integral on the right is finite because the closed
ball is a compact subset of the domain.

\textbf{On the normalisation and the scale.}  The conclusion is
scale-invariant: if $v$ and $(g_\theta)$ satisfy the hypotheses on $B_R(x_0)$
then $\tilde v(y)=v(x_0+Ry)$ and the correspondingly rescaled gradients satisfy
them on $B_1(0)$ with the same $Q$, and the two conclusions correspond under
that change of variables, because the factor $R^{-n}$ is exactly the volume
normalisation.  The inequality is therefore the scale-invariant form of the
estimate ``$\sup_{B_{1/2}}v\le C(n,Q)\bigl(\int_{B_1}v^2\bigr)^{1/2}$'' that the
thesis obtains at $R=1$ from the reverse-H\"older family and then rescales; the
proof below runs at a general $R$ throughout, so no rescaling is needed.  The
supremum is the essential one (p.~6 of the paper), so the conclusion is stated
as an almost-everywhere bound.

\textbf{What the node is for.}  The thesis discharges this step in one sentence
--- ``We can now use Moser's iteration technique as in [GT, chapter 8]'' (p.~17)
--- and the Note for the formalization closes its list of citable results with
``Everything else must be proved'', so Moser's technique must be carried out.
This lemma is that step, isolated from the equation: it mentions no operator $A$,
no solution $u$ and neither of the structure conditions, exactly as \S2 of the
thesis isolates the De Giorgi-type lemmas from the equation.
\end{lemma}
\begin{proof}
Throughout, balls are concentric at $x_0$ and $|E|$ denotes Lebesgue measure;
$|B_\rho|=\omega_n\rho^n$ with $\omega_n=|B_1(0)|$ (Preamble).  Let $\chi>1$ and
$C_S>0$ be the constants furnished by \noderef{SobolevExponentTwoBall} for this
$n$: its two hypotheses before the constants are $2\le n$, which is a hypothesis
of the present lemma, and \noderef{CitedSobolevPoincareBall} in dimension $n$,
which is also a hypothesis of the present lemma and is carried here for exactly
this use.  Both constants depend only on $n$.
Enlarging $C_S$ if necessary we may and do assume $C_S\ge1$, which preserves its
inequality.  For $s\ge1$ and $R/2\le\rho\le\frac34R$ set
\[
Y(s,\rho)=\Bigl(\frac1{|B_R|}\int_{B_\rho}v^{2s}\,dx\Bigr)^{\frac1{2s}},
\qquad
S(\rho)=\operatorname*{ess\,sup}_{B_\rho}v,
\qquad
Y_2=\Bigl(\frac1{|B_R|}\int_{B_R}v^2\,dx\Bigr)^{\frac12}.
\]
These are well defined: $v$ is measurable and nonnegative, and for
$\rho\le\frac34R$ the closed ball $\overline B_\rho$ is a compact subset of
$B_R$, so hypothesis (1) at the exponent $s$ makes $v^s$ square integrable there
and $Y(s,\rho)<\infty$ for every $s>1$; $Y(1,\rho)\le Y_2<\infty$ by the
integrability hypothesis.  We also record once that $S$ is nondecreasing, and
that $v\le S(\rho)$ almost everywhere on $B_\rho$ (Preamble: a function does not
exceed its essential supremum almost everywhere, the bound used in clause (1) of
the item on the limit of normalised $L^q$ means).

\emph{Step 0: the degenerate case.}  If $\int_{B_R}v^2\,dx=0$ then, $v$ being
nonnegative and measurable, $v=0$ almost everywhere on $B_R$ (Preamble: a
nonnegative function with vanishing integral is almost everywhere zero), so
$v\le0$ almost everywhere on $B_{R/2}$ and the conclusion holds for any $C>0$,
both sides being $0$.  Assume from now on $\int_{B_R}v^2\,dx>0$, i.e.\ $Y_2>0$.

\emph{Step 1: the cut-off.}  Let $\rho_{1/8}$ be the standard mollifier of the
Preamble in dimension one at scale $\tfrac18$ --- smooth, nonnegative, supported
in $[-\tfrac18,\tfrac18]$ and of total integral $1$ --- and put
$\psi=\mathbf1_{(-\infty,1/2]}*\rho_{1/8}$.  Then $\psi\in C^\infty(\R)$ with
$\psi'=\mathbf1_{(-\infty,1/2]}*\rho_{1/8}'$ (Preamble: convolution with a smooth
compactly supported kernel, the indicator being locally integrable),
$0\le\psi\le1$,
$\psi=1$ on $(-\infty,\tfrac38]$ and $\psi=0$ on $[\tfrac58,\infty)$ --- for
$t\le\tfrac38$ every $s$ with $|s|\le\tfrac18$ has $t-s\le\tfrac12$, so the
convolution integrates $\rho_{1/8}$ over its whole support, and for
$t\ge\tfrac58$ every such $s$ has $t-s\ge\tfrac12$, so the integrand vanishes
almost everywhere --- and
$|\psi'|\le\|\rho_{1/8}'\|_{L^1(\R)}=:\Lambda<\infty$, an absolute constant,
because $|\mathbf1_{(-\infty,1/2]}|\le1$ and $\rho_{1/8}'$ is continuous with
compact support.  Given radii
$R/2\le r'<r\le\frac34R$ put
\[
\zeta(x)=\psi\Bigl(\frac{|x-x_0|-r'}{r-r'}\Bigr).
\]
Then $0\le\zeta\le1$; $\zeta=1$ on $\overline B_{r'}$, since there the argument
of $\psi$ is $\le0$; and $\zeta$ is smooth, because $x\mapsto|x-x_0|$ is smooth
away from $x_0$ and $\zeta$ is constant on the neighbourhood $B_{r'}$ of $x_0$.

\emph{The support of $\zeta$.}  As everywhere in this tablet,
$\operatorname{supp}$ denotes the \emph{closed} support, $\operatorname{supp}
\zeta=\overline{\{x:\zeta(x)\ne0\}}$, and that is what the clauses this $\zeta$
is fed into bind: hypothesis (2) of the present lemma,
\noderef{WeakGradientSmoothMul} in Step 2 and \noderef{SobolevExponentTwoBall}
in Step 3.  Now $\zeta=0$ outside $\overline B_{r'+\frac58(r-r')}$, because
there the argument of $\psi$ is $\ge\tfrac58$.  That set is \emph{closed}, so it
contains not merely $\{\zeta\ne0\}$ but its closure:
\[
\operatorname{supp}\zeta\ \subseteq\ \overline B_{r'+\frac58(r-r')}\ \subseteq\ B_r ,
\]
the last inclusion because $r'+\tfrac58(r-r')<r$ --- equivalently
$\tfrac58(r-r')<r-r'$, which holds since $r>r'$ --- and a closed ball of radius
strictly less than $r$ about $x_0$ lies in the open ball $B_r$.  In particular
$\operatorname{supp}\zeta$ is a closed subset of the compact
$\overline B_{r'+\frac58(r-r')}$, hence compact, i.e.\ $\zeta$ has compact
support.  Finally
$|\nabla\zeta|=|\psi'|/(r-r')\le\Lambda/(r-r')$, and $\nabla\zeta=0$ at every
point outside $\operatorname{supp}\zeta$.  \emph{This last step is the one place
where the closedness of the support does work, so it is worth saying once}: the
complement of $\operatorname{supp}\zeta$ is open and $\zeta$ vanishes
identically on it, so every derivative of $\zeta$ vanishes there as well.  The
same inference from the bare set $\{\zeta\ne0\}$ would not be available, that
set not being closed in general.  So $\nabla\zeta$ too vanishes outside
$\operatorname{supp}\zeta\subseteq B_r$.

\emph{Step 2: $v^\theta\zeta$ is an admissible test function for the Sobolev
inequality.}  Fix $\theta>1$ and radii $R/2\le r'<r\le\frac34R$, write
$f=v^\theta$, $g=g_\theta$ and $\varphi=f\zeta$ with $\zeta$ as in Step 1.

First, $f\in W^{1,2}(B_r)$ with weak gradient $g$.  Hypothesis (1)
(\noderef{MemW1pLocOn}) gives, on the one hand, that $g$ is the weak gradient of
$f$ on $B_R$ (\noderef{HasWeakGradientOn}) and, on the other, that $f$ and $g$
are square integrable on every compact subset of $B_R$; $\overline B_r$ is such a
subset because $r\le\frac34R<R$, so $f,g\in L^2(\overline B_r)$ and hence
$f,g\in L^2(B_r)$: indeed $|f|^2\mathbf 1_{B_r}\le|f|^2\mathbf
1_{\overline B_r}$ pointwise with the right-hand side integrable, so clause (a)
of the Preamble's domination proposition applies, and likewise for $g$.  For
the weak-gradient identity on the smaller set, we must verify it for the whole
test class that \noderef{HasWeakGradientOn} quantifies over at $B_r$, which is
the triple: $\eta$ of class $C^\infty$ on $\R^n$, $\eta$ of compact support, and
$\operatorname{supp}\eta\subseteq B_r$.  So let $\eta$ be such a test function
and let $i\in\{1,\dots,n\}$.  Since $r\le\frac34R<R$ we have
$B_r\subseteq B_R$ and therefore $\operatorname{supp}\eta\subseteq B_R$, so
$\eta$ also belongs to the test class of \noderef{HasWeakGradientOn} at $B_R$
and the identity there gives
$\int_{B_R}f\,D_i\eta\,dx=-\int_{B_R}g_i\,\eta\,dx$.  Both integrands vanish at
every point outside $\operatorname{supp}\eta$: there $\eta=0$, and
$D_i\eta=0$ by the openness argument recorded in Step 1, the complement of the
closed support being open with $\eta$ identically $0$ on it.  They are integrable on $\overline B_r$, being products of an
$L^2(\overline B_r)$ function with a bounded one.  Since
$\operatorname{supp}\eta\subseteq B_r\subseteq B_R$, each integral over $B_R$
therefore equals the corresponding integral over $B_r$.  Hence
$\int_{B_r}f\,D_i\eta=-\int_{B_r}g_i\eta$ for every $\eta$ in the test class at
$B_r$, which is exactly the assertion that $g$ is the weak gradient of $f$ on
$B_r$ (\noderef{HasWeakGradientOn}), and therefore
$f\in W^{1,2}(B_r)$ with weak gradient $g$ (\noderef{MemW1pOn}).

Now $1\le n$ and $1\le2$; $B_r$ is open; $\zeta$ is of class $C^\infty$, has
compact support and satisfies $\operatorname{supp}\zeta\subseteq B_r$, which are
the three cut-off conditions \noderef{WeakGradientSmoothMul} binds and which
Step 1 established in exactly that form; $|B_r|<\infty$, so
$\mathrm{volume}\,B_r\ne\infty$; and $f\in W^{1,2}(B_r)$ with weak gradient $g$
(\noderef{MemW1pOn}), just shown.  So clause (1) of
\noderef{WeakGradientSmoothMul} applies with $q=2$, $\Omega=B_r$, the pair
$(f,g)$ and $\psi=\zeta$: the product $\varphi=f\zeta$ lies in
$W^{1,2}(B_r)$ with weak gradient
\[
g_\varphi=\zeta\,g+f\,\nabla\zeta .
\]
Finally $\varphi=f\zeta$ vanishes wherever $\zeta$ does, so
$\{\varphi\ne0\}\subseteq\{\zeta\ne0\}$ and hence, taking closures,
$\operatorname{supp}\varphi\subseteq\operatorname{supp}\zeta\subseteq B_r$; and
$\operatorname{supp}\varphi$, a closed subset of the compact
$\operatorname{supp}\zeta$, is compact, so $\varphi$ has compact support.  These
are the two support conditions \noderef{SobolevExponentTwoBall} binds on its
test function.

\emph{Step 3: the reverse-H\"older step.}  Keep $\theta,r',r,\zeta,\varphi$ as in
Step 2.  Apply \noderef{SobolevExponentTwoBall} at the radius $\rho=r$ to the
pair $(\varphi,g_\varphi)$.  Its four remaining hypotheses are: $0<r$, which
holds because $r>r'\ge R/2>0$; $\varphi\in W^{1,2}(B_r)$ with weak gradient
$g_\varphi$ (\noderef{MemW1pOn}), which is Step 2; and the compactness of
$\operatorname{supp}\varphi$ together with
$\operatorname{supp}\varphi\subseteq B_r$, also Step 2.  It gives
\[
\Bigl(\frac1{|B_r|}\int_{B_r}|\varphi|^{2\chi}dx\Bigr)^{\frac1{2\chi}}
\ \le\ C_S\,r\,\Bigl(\frac1{|B_r|}\int_{B_r}|g_\varphi|^2dx\Bigr)^{\frac12}.
\]
We bound the two sides.  On the left, $\zeta=1$ on $B_{r'}$, so $\varphi=v^\theta$
there, and $|\varphi|^{2\chi}\ge0$ elsewhere; together with
$|B_r|^{-1}\ge|B_R|^{-1}$ this gives
\[
\Bigl(\frac1{|B_R|}\int_{B_{r'}}v^{2\theta\chi}dx\Bigr)^{\frac1{2\chi}}
\ \le\
\Bigl(\frac1{|B_r|}\int_{B_r}|\varphi|^{2\chi}dx\Bigr)^{\frac1{2\chi}} .
\]
On the right, $|g_\varphi|^2=|\zeta g+v^\theta\nabla\zeta|^2
\le2\zeta^2|g|^2+2(v^\theta)^2|\nabla\zeta|^2$, by
$|a+b|^2\le(|a|+|b|)^2\le2|a|^2+2|b|^2$ (the triangle inequality and
$2st\le s^2+t^2$); both terms vanish outside
$\operatorname{supp}\zeta\subseteq B_r$ --- the first because $\zeta$ does, the
second because $\nabla\zeta$ does, both by Step 1 --- so their integrals over
$B_r$ and over $B_R$ agree.  Hypothesis (2) of the present lemma may now be
applied at this $\theta$ and this $\zeta$: it quantifies over cut-offs of class
$C^\infty$ with compact support and support inside $B_R$, and
Step 1 gives $\operatorname{supp}\zeta\subseteq B_r\subseteq B_R$.  It gives
\[
\int_{B_r}|g_\varphi|^2dx
\ \le\ 2\int_{B_R}|g|^2\zeta^2dx+2\int_{B_R}(v^\theta)^2|\nabla\zeta|^2dx
\ \le\ (2Q+2)\int_{B_R}(v^\theta)^2|\nabla\zeta|^2dx .
\]
Since $|\nabla\zeta|\le\Lambda/(r-r')$ and $\nabla\zeta$ vanishes outside
$\operatorname{supp}\zeta\subseteq B_r$ (Step 1),
the last integral is at most $\Lambda^2(r-r')^{-2}\int_{B_r}v^{2\theta}dx$.
Finally $|B_r|^{-1}=(R/r)^n|B_R|^{-1}\le2^n|B_R|^{-1}$ because $r\ge R/2$, and
$r\le R$.  Combining,
\[
\Bigl(\frac1{|B_R|}\int_{B_{r'}}v^{2\theta\chi}dx\Bigr)^{\frac1{2\chi}}
\ \le\ D\,\frac R{r-r'}
\Bigl(\frac1{|B_R|}\int_{B_r}v^{2\theta}dx\Bigr)^{\frac12},
\qquad
D=C_S\,2^{\frac n2}(2Q+2)^{\frac12}\Lambda ,
\]
Here $D=D(n,Q)$ depends only on $n$ and $Q$: $C_S$ and $\chi$ depend only on
$n$, and $\Lambda=\|\rho_{1/8}'\|_{L^1(\R)}$ is an absolute constant.  Enlarging
$D$ if necessary we may and do assume $D\ge1$, which preserves the last display:
its right-hand side is nondecreasing in $D$, because $R/(r-r')>0$ and the
remaining factor is nonnegative.  (Only an upper bound on $\Lambda$ was
established in Step 1, so no lower bound on $D$ may be read off its factors; the
enlargement is what makes $D\ge1$ available, and it is used in Step 4 to bound
the partial products by the infinite one.)
Raising both sides of the last display --- both nonnegative and finite --- to the
power $1/\theta$, which is monotone, and recalling the notation of the preamble
to this proof,
\[
Y(\theta\chi,r')\ \le\ \Bigl(D\,\frac R{r-r'}\Bigr)^{\frac1\theta}Y(\theta,r)
\qquad\Bigl(\theta>1,\ \tfrac R2\le r'<r\le\tfrac34R\Bigr).
\tag{1}
\]

\emph{Step 4: the exponent iteration.}  Fix radii $R/2\le\rho'<\rho\le\frac34R$
and set
\[
\theta_j=2\chi^j,\qquad r_j=\rho'+2^{-j}(\rho-\rho')\qquad(j=0,1,2,\dots),
\]
so that $\theta_0=2$, $\theta_{j+1}=\theta_j\chi$, $r_0=\rho$, $r_j$ decreases to
$\rho'$ and $r_j-r_{j+1}=2^{-j-1}(\rho-\rho')$.  Every $\theta_j$ exceeds $1$ and
every $r_j$ lies in $[R/2,\frac34R]$, so $(1)$ applies to each consecutive pair
and gives
\[
Y(\theta_{j+1},r_{j+1})\ \le\ K_j^{\frac1{\theta_j}}\,Y(\theta_j,r_j),
\qquad
K_j=\frac{2^{j+1}DR}{\rho-\rho'} .
\]
All the quantities $Y(\theta_j,r_j)$ are finite, as noted at the outset, so
multiplying these inequalities for $j=0,\dots,m-1$ is legitimate and yields
\[
Y(\theta_m,r_m)\ \le\ \Bigl(\prod_{j=0}^{m-1}K_j^{\frac1{\theta_j}}\Bigr)
Y(2,\rho).
\]
Each $K_j\ge1$, because $D\ge1$ and $R/(\rho-\rho')\ge4$ (the two radii lie in
$[R/2,\frac34R]$, so $\rho-\rho'\le R/4$), so the partial products increase and
are bounded by the infinite one, which we now evaluate.  Writing
$K_j=2^jK_0$ with $K_0=2DR/(\rho-\rho')$,
\[
\sum_{j\ge0}\frac1{\theta_j}=\frac12\sum_{j\ge0}\chi^{-j}=\frac12\frac\chi{\chi-1}=:\sigma_1,
\qquad
\sum_{j\ge0}\frac j{\theta_j}=\frac12\sum_{j\ge0}j\chi^{-j}
=\frac12\frac\chi{(\chi-1)^2}=:\sigma_2 ,
\]
both finite because $\chi>1$ (a geometric and a differentiated geometric series),
and both depending only on $n$.  Hence
\[
\prod_{j\ge0}K_j^{\frac1{\theta_j}}
=K_0^{\sigma_1}2^{\sigma_2}
=\Bigl(\frac{2DR}{\rho-\rho'}\Bigr)^{\sigma_1}2^{\sigma_2},
\]
and therefore, for every $m$,
\[
Y(\theta_m,\rho')\ \le\ Y(\theta_m,r_m)
\ \le\ E\Bigl(\frac R{\rho-\rho'}\Bigr)^{\sigma_1}Y(2,\rho),
\qquad E:=(2D)^{\sigma_1}2^{\sigma_2}=E(n,Q),
\]
the first inequality because $B_{\rho'}\subseteq B_{r_m}$ and the integrand is
nonnegative.  Now let $m\to\infty$.  Since
$0<|B_{\rho'}|<\infty$ and $v$ is measurable,
\[
\Bigl(\frac1{|B_{\rho'}|}\int_{B_{\rho'}}v^{2\theta_m}dx\Bigr)^{\frac1{2\theta_m}}
\ \longrightarrow\ S(\rho')
\]
as $2\theta_m\to\infty$ (Preamble: the essential supremum as the limit of
normalised $L^q$ means, clause (3), applied on $B_{\rho'}$); and
$Y(\theta_m,\rho')$ differs from the left-hand side by the factor
$(|B_{\rho'}|/|B_R|)^{1/(2\theta_m)}$, which tends to $1$ because
$|B_{\rho'}|/|B_R|\in(0,1]$ is a fixed positive number and
$1/(2\theta_m)\to0$.  Hence $Y(\theta_m,\rho')\to S(\rho')$, and
\[
S(\rho')\ \le\ E\Bigl(\frac R{\rho-\rho'}\Bigr)^{\sigma_1}Y(2,\rho)
\qquad\Bigl(\tfrac R2\le\rho'<\rho\le\tfrac34R\Bigr).
\tag{2}
\]
In particular $S(\rho')<\infty$ for every $\rho'<\frac34R$; applied with
$\rho'=\frac58R$ and $\rho=\frac34R$, $(2)$ gives the finite bound
\[
M:=S\bigl(\tfrac58R\bigr)\ \le\ E\,8^{\sigma_1}\,Y\bigl(2,\tfrac34R\bigr)\ <\ \infty .
\tag{3}
\]

\emph{Step 5: from the $L^4$ mean back to the $L^2$ mean.}  Let
$R/2\le\rho\le\frac58R$.  Since $v\le S(\rho)$ almost everywhere on $B_\rho$ and
$v\ge0$, we have $v^4\le S(\rho)^2v^2$ almost everywhere there, so
\[
Y(2,\rho)^4=\frac1{|B_R|}\int_{B_\rho}v^4dx
\ \le\ S(\rho)^2\,\frac1{|B_R|}\int_{B_\rho}v^2dx
\ \le\ S(\rho)^2Y_2^2 ,
\]
that is $Y(2,\rho)\le S(\rho)^{1/2}Y_2^{1/2}$.  Substituting this into $(2)$,
\[
S(\rho')\ \le\ E\Bigl(\frac R{\rho-\rho'}\Bigr)^{\sigma_1}
S(\rho)^{\frac12}\,Y_2^{\frac12}
\qquad\Bigl(\tfrac R2\le\rho'<\rho\le\tfrac58R\Bigr).
\tag{4}
\]

\emph{Step 6: the radius iteration, and the constant.}  Put
\[
\rho_i=\frac R2+\bigl(1-2^{-i}\bigr)\frac R8\qquad(i=0,1,2,\dots),
\]
so $\rho_0=R/2$, the $\rho_i$ increase to $\frac58R$, all lie in
$[\frac R2,\frac58R]$, and
$\rho_{i+1}-\rho_i=2^{-i-1}\frac R8=2^{-i-4}R$, whence
$R/(\rho_{i+1}-\rho_i)=2^{i+4}$.  By $(4)$ applied to the pair
$(\rho_i,\rho_{i+1})$,
\[
S(\rho_i)\ \le\ E\,2^{(i+4)\sigma_1}\,S(\rho_{i+1})^{\frac12}\,Y_2^{\frac12}
\qquad(i\ge0).
\tag{5}
\]
If $S(R/2)=0$ the conclusion is immediate, as in Step 0; so assume
$S(R/2)>0$.  Then $S(\rho_i)\ge S(\rho_0)>0$ for every $i$, and
$S(\rho_i)\le M<\infty$ by $(3)$, so $s_i:=\log S(\rho_i)$ is a finite real
number for every $i$ and $(5)$ reads
\[
s_i\ \le\ \log E+(i+4)\sigma_1\log2+\tfrac12\log Y_2+\tfrac12s_{i+1}.
\]
Iterating this from $i=0$ to $i=m-1$ --- each step multiplies the remaining tail
by $\tfrac12$ --- gives
\[
s_0\ \le\ \sum_{i=0}^{m-1}2^{-i}\Bigl[\log E+(i+4)\sigma_1\log2+\tfrac12\log Y_2\Bigr]
+2^{-m}s_m .
\]
Here $2^{-m}s_m\le2^{-m}\log M\to0$, and the three series converge:
$\sum_{i\ge0}2^{-i}=2$ and $\sum_{i\ge0}(i+4)2^{-i}=2+8=10$.  Letting
$m\to\infty$,
\[
s_0\ \le\ 2\log E+10\,\sigma_1\log2+\log Y_2 ,
\qquad\text{i.e.}\qquad
S\bigl(\tfrac R2\bigr)\ \le\ E^2\,2^{10\sigma_1}\,Y_2 .
\]
Finally $|B_R|=\omega_nR^n$, so
$Y_2=\omega_n^{-1/2}\bigl(R^{-n}\int_{B_R}v^2dx\bigr)^{1/2}$, and with
\[
C:=E^2\,2^{10\sigma_1}\,\omega_n^{-\frac12}>0 ,
\]
which depends only on $n$ and $Q$ --- through $\chi$ and $C_S$, which depend only
on $n$, through the absolute constant $\Lambda$, and through $Q$ itself ---
we obtain
\[
\operatorname*{ess\,sup}_{B_{R/2}}v\ \le\ C\Bigl(R^{-n}\int_{B_R}v^2dx\Bigr)^{\frac12}.
\]
Since $v$ does not exceed its essential supremum almost everywhere (Preamble,
as recorded at the outset of this proof), this is the assertion.
\end{proof}
